\section{Related Literature}
\label{sec:lit}

This paper first relates to the economics of organizations on authority, communication, and knowledge hierarchies. Formal and real authority have long been central to organizational design \citep{AghionTirole1997}. Communication frictions also shape the efficient allocation of decision rights inside organizations \citep{Dessein2002}. More broadly, the internal organization of knowledge is a core determinant of hierarchy and production \citep{Garicano2000}. The present paper builds on this tradition, but adds a dynamic organizational consideration: authority over current tasks influences the firm's future recovery capability when an important execution technology is externally supplied.

The paper is also closely related to recent work on authority allocation between humans and artificial intelligence. \citet{AtheyBryanGans2020} study how decision authority should be allocated between human and artificial decision-makers. My analysis differs in two respects. First, I model AI not as a one-shot recommendation technology, but as a multi-step execution technology, thereby moving beyond the standard treatment of AI as prediction \citep{AgrawalGansGoldfarb2019}. Second, I emphasize that current delegation decisions affect future internal recovery capability. This channel implies that the case against full automation may persist even when AI dominates human handling task by task in normal operations.

More specifically, the paper is intended to speak to a set of questions that have recently become salient in JITE: how technical change interacts with internal control, how governance adapts to changing organizational environments, and how institutional arrangements shape the use of externally supplied technologies. \citet{MarschakWei2024} show how technical change alters decentralization incentives through the interaction with moral hazard. \citet{MukherjeeSaha2024} analyze how delegation interacts with wage bargaining and innovation. \citet{BakerGibbonsMurphy2023} emphasize that organizational design involves not only incentives, but also control and adaptation through the interaction of formal and relational governance. My paper fits naturally into this broader line of inquiry by asking how a new external execution technology changes the boundary between delegated and retained authority inside the firm.

Finally, the paper connects the literature on technical change and delegation to the economics of platforms and institutional dependence. Because the AI capability considered here is supplied by an external platform, its pricing affects not only whether the technology is adopted, but also how authority is allocated within the firm. In that sense, the paper also speaks to the economics of platforms \citep{RochetTirole2003} and to more recent work on platform-side control and self-preferencing \citep{Kim2025}. The generative-AI setting adds an additional layer of dependence and power asymmetry, which has begun to attract attention in adjacent literatures \citep{NzembayieUrbano2026}. Existing work has not provided a comparably simple theory in which dependence on an external agentic-AI platform jointly shapes current authority allocation, retained human recovery capability, and pricing-induced organizational distortion.
